        \btabularx
        \hline
        \multicolumn{4}{|l|}{\bf selectvalue mode attributes} \\
        \hline
        domain     & 3PL & clock   & force the clock domain                    \\     
        direct     & 3PL & boolean & enable direct from test                   \\     
        readonly   & 3PL & boolean & set to read only                          \\     
        alu        & 3P  & boolean & use an ALU if appropriate                 \\        
        dsp        & 3PL & boolean & use a DSP block                           \\
        threestate & 3PL & boolean & use 3-state buffers rather than CLB logic \\
        \hline
        \btabularxend

        {\bf domain} forces the clock domain (both input and output) to
        the specified clock domain. If assigned {\bf null} the current
        clock domain is selected. If the variable already has a clock domain
        as a result of having been assigned or evaluated then that clock domain
        is compared with the intended domain. If the assigned clock domain
        agrees with that already in place then there is no further action,
        otherwise a fatal error occurs.

        {\bf direct} true specifies that if possible the output enable will be
        driven directly from a conditional signal, not through an execution
        thread and WHEN block. This can significantly simplify the logic and
        reduce the probability of failing to meet timing constraints locally.
        Note that when this attribute is in effect the selectvalue is operative
        following configuration loading regardless of any 3PL startup mechanisms
        - see {\bf start()}(\autorefph{sec:ipstart}).
        
        {\bf readonly} true specifies that from now on this variable cannot
        be assigned to. This attribute cannot be assigned false.
        
        {\bf alu} specifies that assignments of expressions containing
        arithmetic operators '+' and '-' as the final or only operation can be
        combined into a single common adder/subtracter using CLBs. If the {\bf alu}
        attribute is false it can override the global directive {\bf ALU} to
        suppress the use of an ALU for an individual static variable.
        
        The {\bf alu} attribute can save logic and incurs no time penalty for the
        operations drawn in to the adder/subtracter, but all other assignments
        suffer an increase in propagation time since assigned values must now
        traverse the shared adder/subtracter. This tradeoff must be considered
        in each case.
        
        The {\bf alu} attribute is only used at completion of immediate execution.
        
        The behaviour of the {\bf dsp} attribute depends on the FPGA -
        \blist
        \item   For Virtex-5 and Virtex-6 it has a similar effect to the
                {\bf alu} attribute in that it draws arithmetic and bitwise logic
                operations in assignments to the variable into a single ALU,
                but using a DSP block rather than CLBs.
        \item   For other FPGAs the attribute has no effect.
        \blend
        
        The {\bf dsp} attribute will not use a Virtex-5 or Virtex-6 DSP block as an
        ALU if the word width exceed 48 bits (multiplier operands are restricted
        to 25 by 18 bits).
        
        Note that for Virtex-5 or Virtex-6 if the {\bf alu} attribute is true but not
        the {\bf dsp} attribute then arithmetic operators '+' and '-' as the final or
        only operation will be combined into a single common adder/subtracter
        using CLBs.
        
        {\bf threestate} true specifies that three state buffers are to be used
        rather than CLB logic. If the device family does not support three state
        drivers then CLB logic will be used regardless of this attribute.
        Devices which have internal three-state buffers are Spartan-II, Spartan
        IIE, Virtex, Virtex-E, virtex-II, Virtex-II Pro, XC9500, XC9500XV and
        XC9500XL.
